

\section{Introduction}
    The field of High Performance Computing (HPC) has been undergoing a transformation towards heterogeneous architectures 
    that are not defined by the sole usage of CPUs. In past decades, the integration of General Purpose Graphic 
    Processing Units (GPGPUs) alongside Central Computing Units (CPUs) have defined the limits of this field's 
    computational power. However in recent years, there has been a rising demand for more advanced systems that have the ability to stimulate
    complex molecular dynamics and combinatorial optimization, shining a light on our current classical computational limitations. 
    Issues such as the curse of dimensionality and the exponential scaling of Hilbert space have redirected this field's focus 
    towards the potential advantages of Quantum Processing Units (QPUs).  QPUs have emerged, not to be the replacement for classical computation, but 
    instead as specialized accelerators that will be integrated into our existing classical systems, as these processing units are capable of handing 
    specific computational tasks that are intractable for classical bits ~\cite{rallis2025interfacingquantumcomputingsystems}. However, quantum technology 
    is a new advancement, one that still needs significant refinement and testing due to their current limitations. 
    This means hybrid frameworks are the current industry goal to work towards in the hopes of achieving new ground within the fields of material science, 
    drug discovery, complex optimization, cryptographic analysis, and many others.  

    Currently, we are in the era of Noisy Intermediate-Scale Quantum (NISQ), where quantum hardware is limited and noisy, 
    unable to harness the true potential of a "quantum advantage" \cite{harrow2017quantumsupremacy}. The most promising path of achieving this 
    is hybrid quantum-classical algorithms ~\cite{rallis2025interfacingquantumcomputingsystems} that combine the computational benefits of each system. 
    Specifically algorithms such as Variational Quantum Eigensolver (VQE), containg a feedback loop that 
    a classical optimizer uses to iteratively refine parameters to be entered into the connected quantum circuit. By offloading a computationally heavy optimization loop to (heterogeneous) classical 
    resources, QPUs prioritze executing circuits and measuring the resulting expectation values, allowing developers  
    to bypass current QPU coherency limitations which leave us with incorrect and decolorant results. However, if long term coherency is achieved, various quantum algorithms could 
    provide us with an exponential speedup over classical methods when determining eigenvalues in large and complex problems ~\cite{peruzzo2014variational}. 
    Classical HPC structures that use MPI, CUDA and GPUs all in dense, parallelized numerical workloads and larger 
    scale simulations, however struggle when the problem dimension grows exponentially alongside 
    the size of the physical system being observed. This issue is present when calculating the arrangement and behavior of electrons in a molecule, but can be addressed 
    by hybrid systems, leaving quantum components to map the wavefunction directly onto qubits while the classical 
    components refine the variational parameters for circuits \cite{rallis2025interfacingquantumcomputingsystems}.
    
    Despite the theoretical potential of hybrid collaboration, this field currently has a significant architectural gap, as  
    there is a lack of specialized middleware that bridges quantum and classical components. Classical HPC frameworks (ones that are optimized 
    for MPI and CUDA) are not designed to communicate directly with quantum control stacks or hardware. This requires quantum researchers and developers to manually handle the complex 
    data serialization and the frequent network latencies that occur in this communication. Existing quantum programming frameworks, such as Qiskit and Cirq, 
    excel with circuit-level QPU control but lack features to effectively integrate with distributed HPC environments, lacking the ability to seamlessly  
    communicate between the two systems. Without the existence of a standardized hybrid middleware, quantum researchers are facing issues such as high tail latencies that are observed when synchronizing
    both HPC nodes and cloud-based QPUs, leading to performances bottlenecks that ruins the possibility of a speedup offered by a
    hybrid system. This issue, combined with the complex task of data serialization for quantum hardware and the potential lack of understanding of the intricate communication  
    between quantum and classical systems, leaves a difficult barrier for researchers to overcome. 

    The proposed project that this paper will undertake is the development of a hybrid quantum-classical software stack that builds the middleware layer needed
    to abstract hardwares and automate the dispatching of tasks. This software stack will design a clear, intuitive API that will abstract the 
    backend of the cloud-based QPU, removing the need for quantum researchers or developers to manually handle the communication complexities in a hybrid system. This hybrid stack 
    will include automated dispatching logic, HPC connectors to the Qiskit framework, and classical kernels 
    that use MPI and CUDA to parallelize the classical optimization subroutines for an added speedup. The stack is benchmarked against VQE workflows 
    on a single host Docker environment alongside MPI/GPU acceleration, aiming to demonstrate measurable progress towards the "quantum advantage," with a 1.5x 
    speedup in the end-to-end wall-clock time when being compared to the baselines of serial, non-distributed and non-optimized classical implementations. Ultimately, 
    the purpose of this work is the design, implementation, and benchmarking an accesible, open-source HPC software framework that achieves a measurable 
    workflow acceleration through integrating QPUs with distributed classical resources. 





